%% kapa: exam=1 round=27 year=2016 period=1 form=A subject=civil_law subject_ko=민법 seq=28 number=28
%% source: Q-Net 감정평가사 기출문제 (cst00309) artlSeq=5169402 file=extracted/1008/5169402/제27회 감정평가사 1차 문제/제27회 감정평가사 1교시문제(A형).pdf
%% conversion: kordoc 4.24.51 -> tools/gen1.py (auto); review=proofread-sonnet; flags=none
\begin{question}{28}
공유에 관한 설명으로 옳지 않은 것은? (다툼이 있으면 판례에 따름)
\begin{choices}
  \item 공유자 사이에 다른 특약이 없는 한 그 지분의 비율로 공유물의 관리비용 기타 의무를 부담한다.
  \item 공유자의 1인이 상속인 없이 사망한 경우, 그 지분은 다른 공유자에게 각 지분의 비율로 귀속된다.
  \item 공유물을 손괴한 자에 대하여 공유자 중 1인은 특별한 사유가 없는 한 공유물에 발생한 손해의 전부를 청구할 수 있다.
  \item 공유토지 위에 건물을 신축하기 위해서는 공유자 전원의 동의가 있어야 한다.
  \item 공유자가 다른 공유자의 지분권을 대외적으로 주장하는 것은 보존행위가 아니다.
\end{choices}
\end{question}
\answer{3}
